\chapter{Collections}
\label{chap:collections}

The \texttt{std::collections} module provides a comprehensive set of data structures optimized for Amiga systems. Each collection is designed with the constraints and opportunities of 68k processors in mind: deterministic memory management, cache-friendly layouts, and zero hidden allocations.

\section{Overview}

Novus collections fall into several categories based on their use cases:

\begin{table}[h]
\centering
\begin{tabular}{lll}
\textbf{Category} & \textbf{Collections} & \textbf{Use Case} \\
\hline
Sequential & \texttt{Vec}, \texttt{VecDeque}, \texttt{SmallVec} & Ordered elements \\
Associative & \texttt{HashMap}, \texttt{HashSet} & Key-value lookup \\
Bit-level & \texttt{BitSet} & Compact boolean storage \\
Streaming & \texttt{RingBuffer}, \texttt{RingBufferPow2} & Fixed-size queues \\
Intrusive & \texttt{IntrusiveList}, \texttt{PriorityList} & AmigaOS interop \\
Arena & \texttt{SlotMap}, \texttt{ObjectPool} & Stable handles \\
\end{tabular}
\caption{Collection categories and their primary use cases}
\end{table}

\subsection{Choosing the Right Collection}

\begin{itemize}
    \item \textbf{Need a growable array?} Use \texttt{Vec<T>}
    \item \textbf{Need a queue or deque?} Use \texttt{VecDeque<T>}
    \item \textbf{Know the typical size?} Use \texttt{SmallVec<T, N>}
    \item \textbf{Need key-value storage?} Use \texttt{HashMap<K, V>}
    \item \textbf{Need unique values?} Use \texttt{HashSet<T>}
    \item \textbf{Need compact flags?} Use \texttt{BitSet}
    \item \textbf{Need a fixed-size buffer?} Use \texttt{RingBuffer<T>}
    \item \textbf{Interfacing with AmigaOS lists?} Use \texttt{IntrusiveList<T>}
    \item \textbf{Need stable handles to objects?} Use \texttt{SlotMap<T>}
    \item \textbf{Need a pool allocator?} Use \texttt{ObjectPool<T>}
\end{itemize}

%%%%%%%%%%%%%%%%%%%%%%%%%%%%%%%%%%%%%%%%%%%%%%%%%%%%%%%%%%%%%%%%%%%%%%%%%%%%%%%
\section{Vec<T> --- Dynamic Array}
\label{sec:vec}
%%%%%%%%%%%%%%%%%%%%%%%%%%%%%%%%%%%%%%%%%%%%%%%%%%%%%%%%%%%%%%%%%%%%%%%%%%%%%%%

\texttt{Vec<T>} is a contiguous, growable array type with heap-allocated contents. It's the most commonly used collection in Novus, similar to Rust's \texttt{Vec} or C++'s \texttt{std::vector}.

\begin{lstlisting}
from std::collections::vec import Vec

fn example() -> Result<(), ExecError> {
    // Create an empty vector
    var numbers = Vec::<u32>::new()

    // Push elements (may allocate)
    numbers.push(1)?
    numbers.push(2)?
    numbers.push(3)?

    // Access by index
    match numbers.get(0) {
        Option::Some(val) => print_u32(val),
        Option::None => {}
    }

    // Pop from the end
    match numbers.pop() {
        Option::Some(val) => print_u32(val),  // 3
        Option::None => {}
    }

    return Result::Ok(())
}
// Memory automatically freed when numbers goes out of scope
\end{lstlisting}

\subsection{Creating Vectors}

\begin{lstlisting}
// Empty vector (zero capacity, no allocation)
var v1 = Vec::<u32>::new()

// Pre-allocated capacity (reduces reallocations)
var v2 = Vec::<u32>::with_capacity(100)?
\end{lstlisting}

\begin{notebox}
\texttt{Vec::new()} performs no allocation. Memory is only allocated on the first \texttt{push()} or \texttt{reserve()} call.
\end{notebox}

\subsection{Core Operations}

\begin{table}[h]
\centering
\begin{tabular}{llll}
\textbf{Method} & \textbf{Returns} & \textbf{Complexity} & \textbf{Description} \\
\hline
\texttt{push(val)} & \texttt{Result<(), ExecError>} & O(1)* & Append to end \\
\texttt{pop()} & \texttt{Option<T>} & O(1) & Remove from end \\
\texttt{get(idx)} & \texttt{Option<T>} & O(1) & Read element \\
\texttt{set(idx, val)} & \texttt{Result<(), ExecError>} & O(1) & Write element \\
\texttt{insert(idx, val)} & \texttt{Result<(), ExecError>} & O(n) & Insert at position \\
\texttt{remove(idx)} & \texttt{Result<T, ExecError>} & O(n) & Remove at position \\
\texttt{len()} & \texttt{u32} & O(1) & Element count \\
\texttt{capacity()} & \texttt{u32} & O(1) & Current capacity \\
\texttt{is\_empty()} & \texttt{bool} & O(1) & Check if empty \\
\texttt{clear()} & \texttt{()} & O(1) & Remove all elements \\
\texttt{reserve(n)} & \texttt{Result<(), ExecError>} & O(n) & Ensure capacity \\
\texttt{shrink\_to\_fit()} & \texttt{Result<(), ExecError>} & O(n) & Reduce capacity \\
\end{tabular}
\caption{\texttt{Vec<T>} method summary. *Amortized O(1), O(n) when growing.}
\end{table}

\subsection{Growth Strategy}

When a \texttt{Vec} needs more capacity, it doubles its current capacity (minimum 4 elements). This amortized doubling strategy provides O(1) push operations on average while minimizing reallocations.

\begin{lstlisting}
var v = Vec::<u32>::new()
// capacity: 0

v.push(1)?   // allocates 4 elements
// capacity: 4

v.push(2)?
v.push(3)?
v.push(4)?
v.push(5)?   // reallocates to 8 elements
// capacity: 8
\end{lstlisting}

\begin{comingfromc}
Unlike C's \texttt{realloc()}, Novus always allocates new memory and copies data. This is because AmigaOS's \texttt{AllocMem}/\texttt{FreeMem} don't support in-place growth. Pre-allocate with \texttt{with\_capacity()} when you know the expected size.
\end{comingfromc}

\subsection{Iteration}

\texttt{Vec<T>} implements the \texttt{Iterable<T>} trait, enabling use in \texttt{for-in} loops:

\begin{lstlisting}
var numbers = Vec::<u32>::new()
numbers.push(10)?
numbers.push(20)?
numbers.push(30)?

for num in numbers {
    print_u32(num)
}
// Prints: 10 20 30
\end{lstlisting}

\subsection{Memory and RAII}

\texttt{Vec<T>} implements \texttt{Drop}, automatically freeing its memory when it goes out of scope:

\begin{lstlisting}
fn process_data() -> Result<(), ExecError> {
    var data = Vec::<u32>::with_capacity(1000)?

    // ... use data ...

    return Result::Ok(())
}  // data.drop() called automatically, memory freed
\end{lstlisting}

%%%%%%%%%%%%%%%%%%%%%%%%%%%%%%%%%%%%%%%%%%%%%%%%%%%%%%%%%%%%%%%%%%%%%%%%%%%%%%%
\section{VecDeque<T> --- Double-Ended Queue}
\label{sec:vecdeque}
%%%%%%%%%%%%%%%%%%%%%%%%%%%%%%%%%%%%%%%%%%%%%%%%%%%%%%%%%%%%%%%%%%%%%%%%%%%%%%%

\texttt{VecDeque<T>} is a double-ended queue implemented as a ring buffer. It provides O(1) insertion and removal at both ends, making it ideal for queues, sliding windows, and work-stealing algorithms.

\begin{lstlisting}
from std::collections::vecdeque import VecDeque

fn message_queue() -> Result<(), ExecError> {
    var queue = VecDeque::<Message>::with_capacity(32)?

    // Enqueue (add to back)
    queue.push_back(msg1)?
    queue.push_back(msg2)?

    // Urgent message to front
    queue.push_front(urgent_msg)?

    // Dequeue (remove from front)
    match queue.pop_front() {
        Option::Some(msg) => process(msg),
        Option::None => {}  // empty
    }

    return Result::Ok(())
}
\end{lstlisting}

\subsection{Ring Buffer Implementation}

Internally, \texttt{VecDeque} uses a circular buffer with \texttt{head} and \texttt{len} indices:

\begin{lstlisting}
// Logical view:    [A, B, C, D]
// Physical memory: [C, D, _, _, A, B]  (head=4, len=4)
//                           ^-- unused space
\end{lstlisting}

This allows both \texttt{push\_front} and \texttt{push\_back} to be O(1) operations without shifting elements.

\subsection{Core Operations}

\begin{table}[h]
\centering
\begin{tabular}{lll}
\textbf{Method} & \textbf{Returns} & \textbf{Description} \\
\hline
\texttt{push\_back(val)} & \texttt{Result<(), ExecError>} & Add to back (enqueue) \\
\texttt{push\_front(val)} & \texttt{Result<(), ExecError>} & Add to front \\
\texttt{pop\_back()} & \texttt{Option<T>} & Remove from back \\
\texttt{pop\_front()} & \texttt{Option<T>} & Remove from front (dequeue) \\
\texttt{front()} & \texttt{Option<T>} & Peek at front \\
\texttt{back()} & \texttt{Option<T>} & Peek at back \\
\texttt{get(idx)} & \texttt{Option<T>} & Access by logical index \\
\texttt{rotate\_left(n)} & \texttt{()} & Rotate elements left \\
\texttt{rotate\_right(n)} & \texttt{()} & Rotate elements right \\
\end{tabular}
\caption{\texttt{VecDeque<T>} method summary}
\end{table}

\subsection{Use Cases}

\begin{itemize}
    \item \textbf{FIFO queues}: \texttt{push\_back()} + \texttt{pop\_front()}
    \item \textbf{LIFO stacks}: \texttt{push\_back()} + \texttt{pop\_back()}
    \item \textbf{Sliding windows}: Maintain a fixed-size window of recent values
    \item \textbf{Priority bypass}: Insert urgent items at front
\end{itemize}

%%%%%%%%%%%%%%%%%%%%%%%%%%%%%%%%%%%%%%%%%%%%%%%%%%%%%%%%%%%%%%%%%%%%%%%%%%%%%%%
\section{SmallVec<T, N> --- Small Vector Optimization}
\label{sec:smallvec}
%%%%%%%%%%%%%%%%%%%%%%%%%%%%%%%%%%%%%%%%%%%%%%%%%%%%%%%%%%%%%%%%%%%%%%%%%%%%%%%

\texttt{SmallVec<T, N>} is a vector that uses the const generic \texttt{N} as its initial capacity. This reduces allocation overhead when you know a typical size.

\begin{lstlisting}
from std::collections::smallvec import SmallVec

// Create a SmallVec with initial capacity of 8
var coords = SmallVec::<Point, 8>::new()

// First push allocates N=8 capacity
coords.push(Point { x: 10, y: 20 })?

// Subsequent pushes use existing capacity
coords.push(Point { x: 30, y: 40 })?
\end{lstlisting}

\begin{notebox}
The current implementation allocates \texttt{N} elements on first push. A future version may store small arrays inline (true small-vector optimization) once Novus supports array types with const generic sizes.
\end{notebox}

\subsection{When to Use SmallVec}

Use \texttt{SmallVec} when:
\begin{itemize}
    \item You know the typical size of your collection
    \item Most instances stay below a certain threshold
    \item You want to avoid the overhead of dynamic resizing
\end{itemize}

\begin{lstlisting}
// A polygon usually has 3-8 vertices
var vertices = SmallVec::<Vertex, 8>::new()

// A function usually has 0-4 parameters
var params = SmallVec::<Param, 4>::new()
\end{lstlisting}

\subsection{Core Operations}

\texttt{SmallVec} provides the same interface as \texttt{Vec}:

\begin{table}[h]
\centering
\begin{tabular}{lll}
\textbf{Method} & \textbf{Returns} & \textbf{Description} \\
\hline
\texttt{push(val)} & \texttt{Result<(), ExecError>} & Append element \\
\texttt{pop()} & \texttt{Option<T>} & Remove last element \\
\texttt{get(idx)} & \texttt{Option<T>} & Read element by index \\
\texttt{get\_mut(idx)} & \texttt{Option<*T>} & Get mutable pointer \\
\texttt{len()} & \texttt{u32} & Element count \\
\texttt{capacity()} & \texttt{u32} & Current capacity \\
\texttt{is\_empty()} & \texttt{bool} & Check if empty \\
\texttt{clear()} & \texttt{()} & Remove all elements \\
\end{tabular}
\caption{\texttt{SmallVec<T, N>} method summary}
\end{table}

%%%%%%%%%%%%%%%%%%%%%%%%%%%%%%%%%%%%%%%%%%%%%%%%%%%%%%%%%%%%%%%%%%%%%%%%%%%%%%%
\section{HashMap<K, V> --- Hash Map}
\label{sec:hashmap}
%%%%%%%%%%%%%%%%%%%%%%%%%%%%%%%%%%%%%%%%%%%%%%%%%%%%%%%%%%%%%%%%%%%%%%%%%%%%%%%

\texttt{HashMap<K, V>} is a hash table with generic keys using open addressing and linear probing. Keys must implement the \texttt{Hash} and \texttt{Eq} traits.

\begin{lstlisting}
from std::collections::hashmap import HashMap

fn symbol_table() -> Result<(), ExecError> {
    var symbols: HashMap<u32, u32> = HashMap::new()?

    // Insert key-value pairs
    symbols.insert(hash("foo"), 100)?
    symbols.insert(hash("bar"), 200)?

    // Lookup by key
    match symbols.get(&hash("foo")) {
        Option::Some(val_ptr) => print_u32(*val_ptr),
        Option::None => print("not found")
    }

    // Check existence
    if symbols.contains_key(&hash("foo")) {
        // key exists
    }

    // Remove
    match symbols.remove(&hash("bar")) {
        Option::Some(old_val) => {},  // was 200
        Option::None => {}
    }

    return Result::Ok(())
}
\end{lstlisting}

\subsection{Implementation Details}

\texttt{HashMap} uses several optimizations for 68k performance:

\begin{itemize}
    \item \textbf{Open addressing}: All entries in a single array (better cache locality)
    \item \textbf{Linear probing}: Simple collision resolution
    \item \textbf{Tombstones}: Lazy deletion for correct probe sequences
    \item \textbf{Power-of-2 capacity}: Fast modulo via bitwise AND
    \item \textbf{70\% load factor}: Automatic resize threshold
    \item \textbf{Cached entry size}: Avoids expensive multiplication
\end{itemize}

\subsection{Core Operations}

\begin{table}[h]
\centering
\begin{tabular}{lll}
\textbf{Method} & \textbf{Returns} & \textbf{Description} \\
\hline
\texttt{insert(k, v)} & \texttt{Result<Option<V>, ExecError>} & Insert/update \\
\texttt{get(\&k)} & \texttt{Option<\&V>} & Lookup by key \\
\texttt{get\_mut(\&k)} & \texttt{Option<\&var V>} & Mutable lookup \\
\texttt{remove(\&k)} & \texttt{Option<V>} & Remove by key \\
\texttt{contains\_key(\&k)} & \texttt{bool} & Check existence \\
\texttt{len()} & \texttt{u32} & Number of pairs \\
\texttt{is\_empty()} & \texttt{bool} & Check if empty \\
\texttt{clear()} & \texttt{()} & Remove all entries \\
\texttt{reserve(n)} & \texttt{Result<(), ExecError>} & Reserve capacity \\
\texttt{compact()} & \texttt{Result<(), ExecError>} & Clear tombstones \\
\texttt{iter()} & \texttt{HashMapIter<K, V>} & Iterator \\
\end{tabular}
\caption{\texttt{HashMap<K, V>} method summary}
\end{table}

\subsection{Iteration}

The iterator returns \texttt{KeyValueRef} structs containing \emph{pointers} to the keys and values:

\begin{lstlisting}
// KeyValueRef structure (returned by iter().next())
struct KeyValueRef<K, V> {
    key: *K,     // Pointer to key
    value: *V,   // Pointer to value
}
\end{lstlisting}

Use pointer dereferencing (\texttt{*}) to access the actual values:

\begin{lstlisting}
var map: HashMap<u32, u32> = HashMap::new()?
map.insert(1, 100)?
map.insert(2, 200)?

var iter = map.iter()
while true {
    match iter.next() {
        Option::Some(kv) => {
            // Note: kv.key and kv.value are pointers
            print_u32(*kv.key)    // Dereference to get key
            print_u32(*kv.value)  // Dereference to get value
        },
        Option::None => break
    }
}
\end{lstlisting}

\begin{warningbox}
Do not modify the map while iterating. This invalidates the iterator and may cause undefined behavior.
\end{warningbox}

\subsection{Key Requirements}

Keys must implement \texttt{Hash} and \texttt{Eq}:

\begin{lstlisting}
struct Point {
    x: i32,
    y: i32,
}

impl Hash for Point {
    fn hash(&self) -> u32 {
        var h: u32 = self.x.hash()
        h = h ^ (self.y.hash() * 2246822507)
        return h
    }
}

impl Eq for Point {
    fn eq(&self, other: &Point) -> bool {
        return self.x == other.x && self.y == other.y
    }
}

// Now Point can be used as a HashMap key
var point_map: HashMap<Point, u32> = HashMap::new()?
\end{lstlisting}

%%%%%%%%%%%%%%%%%%%%%%%%%%%%%%%%%%%%%%%%%%%%%%%%%%%%%%%%%%%%%%%%%%%%%%%%%%%%%%%
\section{HashSet<T> --- Hash Set}
\label{sec:hashset}
%%%%%%%%%%%%%%%%%%%%%%%%%%%%%%%%%%%%%%%%%%%%%%%%%%%%%%%%%%%%%%%%%%%%%%%%%%%%%%%

\texttt{HashSet<T>} is a set of unique values backed by \texttt{HashMap<T, ()>}. It provides O(1) insert, remove, and contains operations.

\begin{lstlisting}
from std::collections::hashset import HashSet

fn unique_ids() -> Result<(), ExecError> {
    var seen: HashSet<u32> = HashSet::new()?

    // Insert returns Ok(true) for new, Ok(false) for existing
    seen.insert(1)?  // Ok(true)
    seen.insert(2)?  // Ok(true)
    seen.insert(1)?  // Ok(false) - already present

    // Check membership
    if seen.contains(&1) {
        // value is in set
    }

    // Remove
    seen.remove(&2)  // returns true

    return Result::Ok(())
}
\end{lstlisting}

\subsection{Set Operations}

\begin{lstlisting}
var a: HashSet<u32> = HashSet::new()?
var b: HashSet<u32> = HashSet::new()?

// Check subset relationship
if a.is_subset(&b) {
    // every element in a is also in b
}

// Check if sets have common elements
if a.intersects(&b) {
    // at least one element in common
}
\end{lstlisting}

%%%%%%%%%%%%%%%%%%%%%%%%%%%%%%%%%%%%%%%%%%%%%%%%%%%%%%%%%%%%%%%%%%%%%%%%%%%%%%%
\section{BitSet --- Compact Bit Storage}
\label{sec:bitset}
%%%%%%%%%%%%%%%%%%%%%%%%%%%%%%%%%%%%%%%%%%%%%%%%%%%%%%%%%%%%%%%%%%%%%%%%%%%%%%%

\texttt{BitSet} provides compact storage for boolean flags using 1 bit per element. This is 32x more memory-efficient than a \texttt{bool} array, essential for resource-constrained Amiga development.

\begin{lstlisting}
from std::collections::bitset import BitSet

fn track_sprites() -> Result<(), ExecError> {
    // 64 bits = 8 bytes (vs 64 bytes for bool array)
    match BitSet::with_capacity(64) {
        Option::Some(var active_sprites) => {
            // Set bits
            active_sprites.set(0)   // sprite 0 is active
            active_sprites.set(5)   // sprite 5 is active

            // Check bits
            if active_sprites.get(0) {
                // sprite 0 is active
            }

            // Count active sprites
            let count = active_sprites.count_ones()  // 2

            // Find first active sprite
            match active_sprites.find_first_set() {
                Option::Some(idx) => {},  // idx == 0
                Option::None => {}
            }
        },
        Option::None => {}
    }

    return Result::Ok(())
}
\end{lstlisting}

\subsection{Bitwise Operations}

\begin{lstlisting}
var a = BitSet::with_capacity(32)?
var b = BitSet::with_capacity(32)?

// Set some bits
a.set(0)
a.set(2)
b.set(1)
b.set(2)

// Bitwise operations (mutate a)
a.and(&b)   // a = a & b (intersection)
a.or(&b)    // a = a | b (union)
a.xor(&b)   // a = a ^ b (symmetric difference)
a.not()     // a = ~a (complement)
\end{lstlisting}

\subsection{Core Operations}

\begin{table}[h]
\centering
\begin{tabular}{lll}
\textbf{Method} & \textbf{Returns} & \textbf{Description} \\
\hline
\texttt{set(idx)} & \texttt{bool} & Set bit to true \\
\texttt{clear(idx)} & \texttt{bool} & Set bit to false \\
\texttt{get(idx)} & \texttt{bool} & Get bit value \\
\texttt{toggle(idx)} & \texttt{bool} & Flip bit \\
\texttt{set\_all()} & \texttt{()} & Set all bits \\
\texttt{clear\_all()} & \texttt{()} & Clear all bits \\
\texttt{count\_ones()} & \texttt{u32} & Population count \\
\texttt{count\_zeros()} & \texttt{u32} & Count cleared bits \\
\texttt{find\_first\_set()} & \texttt{Option<u32>} & First set bit \\
\texttt{find\_first\_clear()} & \texttt{Option<u32>} & First clear bit \\
\texttt{all()} & \texttt{bool} & All bits set? \\
\texttt{any()} & \texttt{bool} & Any bit set? \\
\texttt{none()} & \texttt{bool} & No bits set? \\
\end{tabular}
\caption{\texttt{BitSet} method summary}
\end{table}

\begin{comingfromc}
\texttt{BitSet} uses Brian Kernighan's algorithm for population count and binary search for first-set-bit. On 68020+, these could be replaced with native \texttt{BFFFO} (find first one) instructions for better performance.
\end{comingfromc}

%%%%%%%%%%%%%%%%%%%%%%%%%%%%%%%%%%%%%%%%%%%%%%%%%%%%%%%%%%%%%%%%%%%%%%%%%%%%%%%
\section{RingBuffer<T> --- Fixed-Size Circular Buffer}
\label{sec:ringbuffer}
%%%%%%%%%%%%%%%%%%%%%%%%%%%%%%%%%%%%%%%%%%%%%%%%%%%%%%%%%%%%%%%%%%%%%%%%%%%%%%%

\texttt{RingBuffer<T>} is a fixed-capacity circular buffer ideal for streaming data, audio buffering, and interrupt-safe producer-consumer patterns.

\begin{lstlisting}
from std::collections::ringbuffer import RingBuffer, RingBufferPow2

fn audio_buffer() {
    // Create a 256-sample buffer
    match RingBuffer::<i16>::with_capacity(256) {
        Option::Some(var buffer) => {
            // Producer (e.g., interrupt handler)
            buffer.push_back(sample)  // returns false if full

            // Consumer (e.g., main task)
            match buffer.pop_front() {
                Option::Some(sample) => play(sample),
                Option::None => {}  // buffer empty
            }
        },
        Option::None => {}
    }
}
\end{lstlisting}

\subsection{Power-of-2 Optimization}

For maximum performance, use \texttt{RingBufferPow2<T>} which uses bitwise AND instead of modulo for index wrapping:

\begin{lstlisting}
// Request 100, get 128 (next power of 2)
match RingBufferPow2::<Sample>::with_capacity(100) {
    Option::Some(var buffer) => {
        // buffer.capacity() == 128
        // ~6x faster push/pop on 68000
    },
    Option::None => {}
}
\end{lstlisting}

\begin{tipbox}
\textbf{Performance on 68000}:
\begin{itemize}
    \item \texttt{RingBuffer::push\_back}: $\sim$150 cycles (uses modulo)
    \item \texttt{RingBufferPow2::push\_back}: $\sim$25 cycles (uses bitmask)
\end{itemize}
That's approximately 6x faster! Always prefer \texttt{RingBufferPow2} for performance-critical code like audio streaming and interrupt handlers.
\end{tipbox}

\subsection{Interrupt Safety}

Ring buffers are safe for single-producer/single-consumer patterns:

\begin{itemize}
    \item \textbf{Producer}: Interrupt handler writes via \texttt{push\_back()}
    \item \textbf{Consumer}: Main task reads via \texttt{pop\_front()}
\end{itemize}

\begin{warningbox}
Ring buffers are NOT safe for multiple producers or multiple consumers without external synchronization (Exec \texttt{Forbid}/\texttt{Permit} or semaphores).
\end{warningbox}

\subsection{Amiga Use Cases}

\begin{itemize}
    \item Paula audio sample streaming (16-bit PCM at 28kHz)
    \item Serial port buffering (up to 115200 baud)
    \item Vertical blank event queues (50Hz PAL)
    \item Trackdisk sector buffering
    \item Interrupt-to-task message passing
\end{itemize}

%%%%%%%%%%%%%%%%%%%%%%%%%%%%%%%%%%%%%%%%%%%%%%%%%%%%%%%%%%%%%%%%%%%%%%%%%%%%%%%
\section{IntrusiveList<T> --- AmigaOS-Compatible Lists}
\label{sec:intrusivelist}
%%%%%%%%%%%%%%%%%%%%%%%%%%%%%%%%%%%%%%%%%%%%%%%%%%%%%%%%%%%%%%%%%%%%%%%%%%%%%%%

\texttt{IntrusiveList<T>} is an intrusive doubly-linked list that matches AmigaOS Exec's \texttt{MinList}/\texttt{MinNode} pattern exactly. It uses native Exec list functions for maximum compatibility.

\begin{lstlisting}
from std::collections::intrusive_list import IntrusiveList, init_min_node
from amiga::raw::structs import MinNode

struct MyData {
    node: MinNode,  // MUST be first field
    value: u32,
}

fn use_list() {
    var list = IntrusiveList::<MyData>::new()
    list.init()  // MUST call after assignment

    // Create node (caller manages lifetime)
    var item = MyData {
        node: MinNode { mln_Succ: null, mln_Pred: null },
        value: 42
    }

    // Add to list
    list.add_tail(&item)

    // Iterate
    match list.first() {
        Option::Some(ptr) => {
            var current = ptr
            while true {
                print_u32((*current).value)
                match list.next(current) {
                    Option::Some(next) => current = next,
                    Option::None => break
                }
            }
        },
        Option::None => {}
    }

    // Remove before freeing
    list.remove(&item)
}
\end{lstlisting}

\begin{warningbox}
You \textbf{MUST} call \texttt{init()} immediately after creating an \texttt{IntrusiveList}. The list contains self-referential pointers that are invalid until \texttt{init()} is called. Using an uninitialized list will corrupt memory or crash.
\end{warningbox}

\subsection{Key Constraints}

\begin{enumerate}
    \item \textbf{Node must be first field}: The \texttt{MinNode} must be the first field of your struct
    \item \textbf{Call init() after assignment}: The list has self-referential pointers that must be set up after the list is at its final location
    \item \textbf{Caller manages node lifetime}: The list does NOT own the nodes
    \item \textbf{One list per node}: A node can only be in one list at a time
\end{enumerate}

\subsection{PriorityList<T>}

\texttt{PriorityList<T>} is similar but uses full \texttt{Node} structures with priority support:

\begin{lstlisting}
from std::collections::intrusive_list import PriorityList, init_node
from amiga::raw::structs import Node

struct Task {
    node: Node,  // Full node with priority
    data: u32,
}

fn priority_queue() {
    var tasks = PriorityList::<Task>::new()
    tasks.init()

    var high = Task { node: @zeroed(Node), data: 1 }
    var low = Task { node: @zeroed(Node), data: 2 }

    // Set priorities
    init_node(&high.node, 10, null)  // priority 10
    init_node(&low.node, 1, null)    // priority 1

    // Insert in priority order (higher = closer to head)
    tasks.enqueue(&high)
    tasks.enqueue(&low)

    // First will be the high-priority task
    match tasks.first() {
        Option::Some(ptr) => {
            // (*ptr).data == 1 (high priority)
        },
        Option::None => {}
    }
}
\end{lstlisting}

%%%%%%%%%%%%%%%%%%%%%%%%%%%%%%%%%%%%%%%%%%%%%%%%%%%%%%%%%%%%%%%%%%%%%%%%%%%%%%%
\section{SlotMap<T> --- Generational Arena}
\label{sec:slotmap}
%%%%%%%%%%%%%%%%%%%%%%%%%%%%%%%%%%%%%%%%%%%%%%%%%%%%%%%%%%%%%%%%%%%%%%%%%%%%%%%

\texttt{SlotMap<T>} is a generational index arena that provides stable handles to objects. It detects use-after-free by tracking a generation counter for each slot.

\begin{lstlisting}
from std::collections::slotmap import SlotMap, SlotKey

struct Entity {
    x: i32,
    y: i32,
    health: u32,
}

fn game_entities() {
    match SlotMap::<Entity>::new() {
        Option::Some(var entities) => {
            // Insert returns a SlotKey handle
            match entities.insert(Entity { x: 100, y: 50, health: 100 }) {
                Option::Some(player_key) => {
                    // Access via handle
                    match entities.get(player_key) {
                        Option::Some(entity) => {
                            print_i32((*entity).x)
                        },
                        Option::None => {}
                    }

                    // Remove the entity
                    entities.remove(player_key)

                    // Old key is now invalid!
                    match entities.get(player_key) {
                        Option::Some(_) => {},  // Never reached
                        Option::None => {}      // Key is stale
                    }
                },
                Option::None => {}
            }
        },
        Option::None => {}
    }
}
\end{lstlisting}

\subsection{How Generations Work}

Each slot has a generation counter. When an object is removed, the generation is incremented. Old keys have the old generation and fail the check:

\begin{lstlisting}
// Slot 0, Generation 0 -> player inserted
let player_key = SlotKey { index: 0, generation: 0 }

// Slot 0, Generation 1 -> player removed
entities.remove(player_key)

// Slot 0 reused for enemy
let enemy_key = SlotKey { index: 0, generation: 1 }

// Old player_key fails - generation mismatch!
entities.get(player_key)  // Returns None
\end{lstlisting}

\subsection{SlotKey}

\texttt{SlotKey} contains both index and generation:

\begin{lstlisting}
struct SlotKey {
    index: u32,
    generation: u32,
}

// Pack into single u32 (16 bits each)
let packed: u32 = key.pack()

// Unpack
let key = SlotKey::unpack(packed)

// Sentinel/null key
let null_key = SlotKey::sentinel()
if key.is_sentinel() { /* invalid */ }
\end{lstlisting}

\subsection{Use Cases}

\begin{itemize}
    \item \textbf{ECS game entities}: Players, enemies, projectiles, pickups
    \item \textbf{Resource handles}: Textures, sounds, fonts
    \item \textbf{Window/Screen handles}: Multiple windows with stable IDs
    \item \textbf{Network object sync}: Stable references across frames
    \item \textbf{Undo/redo systems}: References that survive modifications
\end{itemize}

\subsection{SlotMapFixed<T>}

\texttt{SlotMapFixed<T>} is a variant with fixed capacity (no growth). It uses separate arrays for generations, occupancy flags, and values, providing better cache behavior during iteration:

\begin{lstlisting}
from std::collections::slotmap import SlotMapFixed

// Fixed capacity - never grows, but more cache-friendly
match SlotMapFixed::<Entity>::with_capacity(1024) {
    Option::Some(var entities) => {
        // Same API as SlotMap
        match entities.insert(player) {
            Option::Some(key) => { /* ... */ },
            Option::None => { /* capacity full */ }
        }
    },
    Option::None => {}
}
\end{lstlisting}

\begin{notebox}
Use \texttt{SlotMap<T>} when you need dynamic growth. Use \texttt{SlotMapFixed<T>} when you know the maximum count upfront and want better iteration performance.
\end{notebox}

%%%%%%%%%%%%%%%%%%%%%%%%%%%%%%%%%%%%%%%%%%%%%%%%%%%%%%%%%%%%%%%%%%%%%%%%%%%%%%%
\section{ObjectPool<T> --- Pool Allocator}
\label{sec:objectpool}
%%%%%%%%%%%%%%%%%%%%%%%%%%%%%%%%%%%%%%%%%%%%%%%%%%%%%%%%%%%%%%%%%%%%%%%%%%%%%%%

\texttt{ObjectPool<T>} is a slab allocator for fixed-size objects with O(1) allocation and deallocation. After initialization, no system calls are needed.

\begin{lstlisting}
from std::collections::freelist import ObjectPool

struct Bullet {
    x: i32,
    y: i32,
    dx: i32,
    dy: i32,
}

fn bullet_system() {
    // Pre-allocate pool of 128 bullets
    match ObjectPool::<Bullet>::with_capacity(128) {
        Option::Some(var pool) => {
            // Allocate a bullet (O(1), no syscall)
            match pool.alloc() {
                Option::Some(bullet) => {
                    (*bullet).x = 100
                    (*bullet).y = 50
                    (*bullet).dx = 5
                    (*bullet).dy = 0

                    // ... use bullet ...

                    // Return to pool (O(1))
                    pool.free(bullet)
                },
                Option::None => {}  // pool exhausted
            }
        },
        Option::None => {}
    }
}
\end{lstlisting}

\subsection{Interrupt Safety}

\texttt{ObjectPool} is safe for single-producer/single-consumer patterns:
\begin{itemize}
    \item \textbf{Producer}: Interrupt handler allocates via \texttt{alloc()}
    \item \textbf{Consumer}: Main task frees via \texttt{free()}
\end{itemize}

\begin{warningbox}
\textbf{Safety requirements}:
\begin{itemize}
    \item Do NOT free objects not allocated from this pool
    \item Do NOT double-free objects
    \item Do NOT use objects after freeing them
\end{itemize}
\end{warningbox}

\subsection{ObjectPoolFixed<T>}

\texttt{ObjectPoolFixed<T>} uses an index-stack approach instead of a linked free list. This provides additional capabilities:

\begin{itemize}
    \item \textbf{Allocation tracking}: Knows exactly which slots are in use
    \item \textbf{Proper Drop support}: Calls \texttt{Drop} on all allocated objects when pool is dropped or reset
    \item \textbf{Index-based iteration}: Can iterate over allocated objects by index
    \item \textbf{Free validation}: Detects double-free and invalid free operations
\end{itemize}

\begin{lstlisting}
from std::collections::freelist import ObjectPoolFixed

match ObjectPoolFixed::<Enemy>::with_capacity(32) {
    Option::Some(var pool) => {
        // Allocate
        match pool.alloc() {
            Option::Some(enemy) => {
                // Initialize...
            },
            Option::None => {}
        }

        // Iterate over all allocated enemies
        var i: u32 = 0
        while i < pool.capacity() {
            if pool.is_allocated(i) {
                match pool.get(i) {
                    Option::Some(enemy) => update(enemy),
                    Option::None => {}
                }
            }
            i = i + 1
        }

        // Reset calls Drop on all allocated objects
        pool.reset()
    },
    Option::None => {}
}
\end{lstlisting}

\begin{notebox}
Use \texttt{ObjectPool<T>} for maximum speed when you don't need iteration or Drop support. Use \texttt{ObjectPoolFixed<T>} when you need to iterate over allocated objects or when \texttt{T} owns resources that must be cleaned up.
\end{notebox}

%%%%%%%%%%%%%%%%%%%%%%%%%%%%%%%%%%%%%%%%%%%%%%%%%%%%%%%%%%%%%%%%%%%%%%%%%%%%%%%
\section{Performance Summary}
\label{sec:collection-performance}
%%%%%%%%%%%%%%%%%%%%%%%%%%%%%%%%%%%%%%%%%%%%%%%%%%%%%%%%%%%%%%%%%%%%%%%%%%%%%%%

\begin{table}[h]
\centering
\small
\begin{tabular}{lcccccc}
\textbf{Collection} & \textbf{Insert} & \textbf{Remove} & \textbf{Access} & \textbf{Search} & \textbf{Memory} \\
\hline
\texttt{Vec<T>} & O(1)* & O(n) & O(1) & O(n) & Contiguous \\
\texttt{VecDeque<T>} & O(1) & O(1)** & O(1) & O(n) & Ring buffer \\
\texttt{HashMap<K,V>} & O(1)* & O(1) & O(1) & O(1) & Hash table \\
\texttt{HashSet<T>} & O(1)* & O(1) & --- & O(1) & Hash table \\
\texttt{BitSet} & O(1) & O(1) & O(1) & O(n) & 1 bit/element \\
\texttt{RingBuffer<T>} & O(1) & O(1) & O(1) & O(n) & Fixed size \\
\texttt{IntrusiveList<T>} & O(1) & O(1) & O(n) & O(n) & No allocation \\
\texttt{SlotMap<T>} & O(1)* & O(1) & O(1) & --- & Arena \\
\texttt{ObjectPool<T>} & O(1) & O(1) & --- & --- & Fixed size \\
\end{tabular}
\caption{Collection complexity summary. *Amortized. **Front/back only.}
\end{table}

\begin{comingfromc}
Unlike C's standard library, all Novus collections:
\begin{itemize}
    \item Use AmigaOS \texttt{AllocMem}/\texttt{FreeMem} for memory management
    \item Are properly aligned for 68k DMA and cache efficiency
    \item Implement RAII via the \texttt{Drop} trait
    \item Return \texttt{Result} or \texttt{Option} instead of null pointers
\end{itemize}
\end{comingfromc}
